\documentclass[11pt]{article}

\usepackage[a4paper,margin=1in]{geometry}
\usepackage{amsmath,amssymb,amsthm,mathtools}
\usepackage[T1]{fontenc}
\usepackage{lmodern}
\usepackage{microtype}
\usepackage{hyperref}

\newtheorem{theorem}{Theorem}[section]
\newtheorem{lemma}[theorem]{Lemma}
\newtheorem{proposition}[theorem]{Proposition}
\newtheorem{corollary}[theorem]{Corollary}
\newtheorem{remark}[theorem]{Remark}
\newtheorem{definition}[theorem]{Definition}
\newtheorem{hypothesis}[theorem]{Hypothesis}

\DeclareMathOperator{\ind}{ind}
\DeclareMathOperator{\diag}{diag}
\DeclareMathOperator{\tridiag}{tridiag}

\newcommand{\R}{\mathbb{R}}
\newcommand{\C}{\mathbb{C}}

\title{Algebraic Reductions and a Conditional Topological Argument\\
for the Clean Hatano--Nelson Family with Open Boundary Conditions}
\author{}
\date{}

\begin{document}
\maketitle

\section{Setup}

Fix \(n\ge 2\), \(0<a<1\), and consider the real nonnormal tridiagonal Toeplitz matrix
\[
\mathbf A_n(a):=\tridiag(a,0,1)
=
\begin{bmatrix}
0 & 1 \\
a & 0 & \ddots \\
& \ddots & \ddots & \ddots \\
&& \ddots & 0 & 1 \\
&&& a & 0
\end{bmatrix}\in \R^{n\times n}.
\]

For \(\varepsilon>0\), its \(\varepsilon\)-pseudospectrum is
\[
\sigma_\varepsilon(\mathbf A_n(a))
:=
\left\{
z\in\C:\;
\|(z\mathbf I_n-\mathbf A_n(a))^{-1}\|_2>\varepsilon^{-1}
\right\},
\]
with the usual convention that the resolvent norm is \(+\infty\) on the spectrum.
Equivalently,
\[
\sigma_\varepsilon(\mathbf A_n(a))
=
\left\{
z\in\C:\;
s_{\min}(z\mathbf I_n-\mathbf A_n(a))<\varepsilon
\right\}.
\]

\begin{lemma}[Real spectrum]\label{lem:real-spectrum}
Let
\[
\mathbf D_n:=\diag(1,q,q^2,\dots,q^{n-1}),
\qquad
q:=\sqrt a.
\]
Then
\[
\mathbf D_n^{-1}\mathbf A_n(a)\mathbf D_n
=
\sqrt a\,\tridiag(1,0,1).
\]
Consequently,
\[
\sigma(\mathbf A_n(a))
=
\left\{
2\sqrt a \cos\frac{j\pi}{n+1}:j=1,\dots,n
\right\}
\subset \R.
\]
\end{lemma}

\begin{proof}
A direct computation gives
\[
(\mathbf D_n^{-1}\mathbf A_n(a)\mathbf D_n)_{j,j+1}=q=\sqrt a,
\qquad
(\mathbf D_n^{-1}\mathbf A_n(a)\mathbf D_n)_{j+1,j}=aq^{-1}=\sqrt a.
\]
Thus \(\mathbf D_n^{-1}\mathbf A_n(a)\mathbf D_n=\sqrt a\,\tridiag(1,0,1)\), and the spectral formula is the classical one for the symmetric tridiagonal Toeplitz matrix.
\end{proof}

\begin{remark}[Two endpoint models]\label{rem:endpoints}
The continuous endpoint families \(a=1\) and \(a=0\) are especially transparent.

\begin{enumerate}
\item If \(a=1\), then \(\mathbf A_n(1)=\tridiag(1,0,1)\) is symmetric, hence normal, and
\[
\sigma_\varepsilon(\mathbf A_n(1))
=
\bigcup_{j=1}^n
D\!\left(2\cos\frac{j\pi}{n+1},\varepsilon\right).
\]
Therefore every connected component is simply connected.

\item If \(a=0\), then \(\mathbf A_n(0)\) is the nilpotent Jordan block with superdiagonal \(1\). For every \(\theta\in\R\),
\[
\mathbf U_\theta^*\,\mathbf A_n(0)\,\mathbf U_\theta
=
e^{i\theta}\mathbf A_n(0),
\qquad
\mathbf U_\theta:=\diag(1,e^{i\theta},\dots,e^{i(n-1)\theta}),
\]
so \(\sigma_\varepsilon(\mathbf A_n(0))\) is rotationally invariant. It is standard that every bounded connected component of a matrix pseudospectrum contains an eigenvalue; see, for example, \cite{TrefethenEmbree2005}. Since \(0\) is the only eigenvalue of \(\mathbf A_n(0)\), the set \(\sigma_\varepsilon(\mathbf A_n(0))\) has exactly one connected component. Rotational invariance then forces it to be a disc centered at \(0\), hence simply connected.
\end{enumerate}

These two limit cases motivate the low-\(a\) and near-normal regimes used in the final topological argument.
\end{remark}

\begin{proposition}[Low-\(a\) star-shaped regime]\label{prop:low-a}
Fix \(n\ge 2\) and \(\varepsilon>0\). Then there exists
\[
a_-=a_-(n,\varepsilon)>0
\]
such that for every
\[
0\le a<a_-,
\]
the set \(\sigma_\varepsilon(\mathbf A_n(a))\) is a single star-shaped Jordan domain.
In particular, the low-\(a\) part later isolated in
Hypothesis~\ref{hyp:standing-input} is not an external input.
\end{proposition}

\begin{proof}
Set
\[
\mathbf J_n:=\mathbf A_n(0).
\]
By Remark~\ref{rem:endpoints}, the set \(\sigma_\varepsilon(\mathbf J_n)\) is a disc
\[
\sigma_\varepsilon(\mathbf J_n)=D(0,\rho)
\]
for a unique radius \(\rho=\rho_{n,\varepsilon}>0\).

We first show that
\[
f_0(r):=s_{\min}(r\mathbf I_n-\mathbf J_n),\qquad r>0,
\]
is strictly increasing. For \(r>0\),
\[
(r\mathbf I_n-\mathbf J_n)^{-1}
=
\sum_{k=0}^{n-1} r^{-k-1}\mathbf J_n^k
=: \mathbf R_r.
\]
The matrix \(\mathbf R_r\) is upper triangular with positive entries, and every nonzero entry is a strictly decreasing function of \(r\). Hence the Hermitian matrix
\[
\mathbf B_r:=\mathbf R_r^*\mathbf R_r
\]
has strictly positive entries, and every entry of \(\mathbf B_r\) is strictly decreasing in \(r\). By the Perron--Frobenius theorem, its largest eigenvalue
\[
\mu(r):=\rho(\mathbf B_r)=\|\mathbf R_r\|_2^2
\]
is simple and has a strictly positive unit eigenvector \(\mathbf w(r)\). Therefore
\[
\mu'(r)=\mathbf w(r)^{\mathsf T}\mathbf B_r'\mathbf w(r)<0
\]
because every entry of \(\mathbf B_r'\) is strictly negative. Since
\[
f_0(r)=\mu(r)^{-1/2},
\]
it follows that \(f_0\) is real-analytic and strictly increasing on \((0,\infty)\). In particular,
\[
f_0(\rho)=\varepsilon.
\]

Next, by the rotational covariance
\[
\mathbf U_\theta^*\,\mathbf J_n\,\mathbf U_\theta=e^{i\theta}\mathbf J_n,
\]
we have
\[
s_{\min}(re^{i\theta}\mathbf I_n-\mathbf J_n)=f_0(r)
\qquad (\theta\in\R).
\]
Moreover, the smallest singular value is simple on the circle \(|z|=\rho\). Indeed,
\[
s_{\min}(r\mathbf I_n-\mathbf J_n)=\|\mathbf R_r\|_2^{-1},
\]
and \(\|\mathbf R_r\|_2^2=\mu(r)\) is simple.

Choose \(\eta>0\) so small that
\[
0<\rho-2\eta<\rho+2\eta<2+\varepsilon,
\]
and
\[
f_0'(r)\ge 2m>0
\qquad \text{for all } r\in[\rho-2\eta,\rho+2\eta]
\]
for some \(m>0\). Since the smallest singular value of
\(re^{i\theta}\mathbf I_n-\mathbf J_n\) is simple on the compact annulus
\[
\rho-2\eta\le r\le \rho+2\eta,
\]
there exists \(\gamma>0\) such that
\[
s_2(re^{i\theta}\mathbf I_n-\mathbf J_n)-s_{\min}(re^{i\theta}\mathbf I_n-\mathbf J_n)\ge 4\gamma
\]
for all \(\theta\in\R\) and all \(r\in[\rho-2\eta,\rho+2\eta]\), where \(s_2\) denotes the second-smallest singular value.

Now
\[
\mathbf A_n(a)=\mathbf J_n+a\mathbf L_n,
\]
where \(\mathbf L_n\) is the lower shift and \(\|\mathbf L_n\|_2=1\). Hence
\[
\|\mathbf A_n(a)-\mathbf J_n\|_2=a.
\]
By Weyl's inequality for singular values, if \(0\le a<\gamma\), then on the annulus
\[
\rho-2\eta\le r\le \rho+2\eta
\]
the smallest singular value of
\[
re^{i\theta}\mathbf I_n-\mathbf A_n(a)
\]
remains simple. Therefore, by standard perturbation theory for simple singular values, there is \(a_1>0\) such that for \(0\le a<a_1\) the function
\[
(r,\theta,a)\longmapsto s_{\min}(re^{i\theta}\mathbf I_n-\mathbf A_n(a))
\]
is real-analytic on
\[
[\rho-2\eta,\rho+2\eta]\times \R \times [0,a_1),
\]
and, after shrinking \(a_1\) if necessary,
\[
\partial_r s_{\min}(re^{i\theta}\mathbf I_n-\mathbf A_n(a))\ge m
\]
throughout that set.

Choose \(c>0\) such that
\[
f_0(r)\le \varepsilon-3c
\quad (0\le r\le \rho-\eta),
\qquad
f_0(r)\ge \varepsilon+3c
\quad (\rho+\eta\le r\le 2+\varepsilon).
\]
Again by Weyl's inequality, if \(0\le a<c\), then
\[
s_{\min}(re^{i\theta}\mathbf I_n-\mathbf A_n(a))\le \varepsilon-2c
\quad (0\le r\le \rho-\eta),
\]
and
\[
s_{\min}(re^{i\theta}\mathbf I_n-\mathbf A_n(a))\ge \varepsilon+2c
\quad (\rho+\eta\le r\le 2+\varepsilon).
\]
Since \(\|\mathbf A_n(a)\|_2\le 2\), every point of \(\sigma_\varepsilon(\mathbf A_n(a))\) satisfies \(|z|<2+\varepsilon\). Thus every boundary point lies in the annulus
\[
\rho-\eta<|z|<\rho+\eta.
\]

Let
\[
a_-:=\min\{a_1,\gamma,c\}.
\]
Fix \(a\in[0,a_-)\). For each \(\theta\in\R\), the function
\[
r\longmapsto s_{\min}(re^{i\theta}\mathbf I_n-\mathbf A_n(a))
\]
is strictly increasing on \([\rho-\eta,\rho+\eta]\), is \(<\varepsilon\) at \(r=\rho-\eta\), and is \(>\varepsilon\) at \(r=\rho+\eta\). Hence there exists a unique radius
\[
r(\theta,a)\in(\rho-\eta,\rho+\eta)
\]
such that
\[
s_{\min}(r(\theta,a)e^{i\theta}\mathbf I_n-\mathbf A_n(a))=\varepsilon.
\]
Because the boundary has no points outside that annulus, every ray from the origin meets
\(\partial\sigma_\varepsilon(\mathbf A_n(a))\) in exactly one point.

Therefore
\[
\sigma_\varepsilon(\mathbf A_n(a))
=
\{re^{i\theta}:\ 0\le r<r(\theta,a)\},
\]
with \(r(\theta,a)\) continuous and \(2\pi\)-periodic in \(\theta\). Hence
\(\sigma_\varepsilon(\mathbf A_n(a))\) is a star-shaped Jordan domain. In particular it is connected.
\end{proof}

\begin{definition}[Noncritical normal-endpoint radius]\label{def:noncritical}
Let
\[
\lambda_j^{(1)}:=2\cos\frac{j\pi}{n+1},
\qquad j=1,\dots,n.
\]
We say that \(\varepsilon>0\) is \emph{noncritical at the normal endpoint} if
\[
2\varepsilon \neq \bigl|\lambda_j^{(1)}-\lambda_k^{(1)}\bigr|
\qquad\text{for all } j\neq k.
\]
Equivalently, no two circles
\[
\partial D(\lambda_j^{(1)},\varepsilon)
\]
and
\[
\partial D(\lambda_k^{(1)},\varepsilon)
\]
are tangent.
\end{definition}

\begin{proposition}[Near-normal regime under the noncriticality condition]\label{prop:near-normal}
Fix \(n\ge 2\) and \(\varepsilon>0\), and assume that \(\varepsilon\) is noncritical at the normal endpoint in the sense of Definition~\ref{def:noncritical}. Then there exists
\[
\widehat a_{n,\varepsilon}\in(0,1)
\]
such that for every
\[
\widehat a_{n,\varepsilon}<a\le 1,
\]
the sets
\[
\sigma_\varepsilon(\mathbf A_n(a))
\quad\text{and}\quad
\sigma_\varepsilon(\mathbf A_n(1))
\]
are semialgebraically homeomorphic. In particular, every connected component of
\[
\sigma_\varepsilon(\mathbf A_n(a))
\]
is simply connected for \(\widehat a_{n,\varepsilon}<a\le 1\).
\end{proposition}

\begin{proof}[Proof sketch]
Set
\[
P_{a,\varepsilon}:=\sigma_\varepsilon(\mathbf A_n(a)).
\]
By Lemma~\ref{lem:real-spectrum},
\[
\sigma(\mathbf A_n(1))
=
\left\{
\lambda_j^{(1)}:j=1,\dots,n
\right\},
\qquad
\lambda_j^{(1)}=2\cos\frac{j\pi}{n+1}.
\]
Since \(\mathbf A_n(1)\) is normal,
\[
P_{1,\varepsilon}
=
\bigcup_{j=1}^n D(\lambda_j^{(1)},\varepsilon).
\]

By Definition~\ref{def:noncritical}, no two boundary circles of this union are tangent. Hence every point of \(\partial P_{1,\varepsilon}\) is of one of the following two types:

\begin{enumerate}
\item it lies on exactly one circle \(\partial D(\lambda_j^{(1)},\varepsilon)\), and is therefore a smooth boundary point;
\item it is an intersection point of exactly two circles, and these two circles meet transversely.
\end{enumerate}

No boundary point lies on three circles, since the centers are distinct and collinear.

Fix \(a_0<1\). By Lemma~\ref{lem:semialg}, the set
\[
\mathcal X
:=
\{(x,y,a)\in \R^2\times[a_0,1]:x+iy\in P_{a,\varepsilon}\}
\]
is bounded and semialgebraic. Hence Hardt's semialgebraic triviality theorem
\cite[Ch.~9]{BCR} implies that the projection
\[
\pi:\mathcal X\to[a_0,1],\qquad (x,y,a)\mapsto a,
\]
has only finitely many exceptional parameters. It is therefore enough to show that \(a=1\) is not exceptional.

Let \(p\in\partial P_{1,\varepsilon}\).

If \(p\) is a smooth point belonging to exactly one circle
\[
\partial D(\lambda_j^{(1)},\varepsilon),
\]
choose a small neighborhood \(U\) of \(p\) avoiding all other circles. In \(U\), the active singular-value branch of
\[
z\mathbf I_n-\mathbf A_n(1)
\]
is the simple function
\[
z\longmapsto |z-\lambda_j^{(1)}|,
\]
and its gradient is nonzero on the level set
\[
|z-\lambda_j^{(1)}|=\varepsilon.
\]
Since singular values depend continuously on the matrix entries, and simple singular-value branches depend smoothly on parameters, the implicit-function theorem shows that for \(a\) sufficiently close to \(1\), the local boundary
\[
U\cap\partial P_{a,\varepsilon}
\]
is a single embedded arc, and the local pair
\[
(U,U\cap P_{a,\varepsilon})
\]
is homeomorphic to
\[
(U,U\cap P_{1,\varepsilon}).
\]

If \(p\) is a transverse intersection point of exactly two circles
\[
\partial D(\lambda_j^{(1)},\varepsilon)
\quad\text{and}\quad
\partial D(\lambda_k^{(1)},\varepsilon),
\]
choose \(U\) so small that no other circle meets \(U\). Then
\[
U\cap P_{1,\varepsilon}
=
U\cap\bigl(D(\lambda_j^{(1)},\varepsilon)\cup D(\lambda_k^{(1)},\varepsilon)\bigr),
\]
and the two boundary arcs meet transversely. The two relevant singular values of
\[
z\mathbf I_n-\mathbf A_n(1)
\]
are
\[
|z-\lambda_j^{(1)}|
\quad\text{and}\quad
|z-\lambda_k^{(1)}|,
\]
and, on \(U\), they are separated from all remaining singular values by a positive gap. Standard perturbation theory for isolated singular-value clusters therefore shows that for \(a\) sufficiently close to \(1\), the local boundary in \(U\) is a small perturbation of this transverse corner configuration. In particular, the local pair
\[
(U,U\cap P_{a,\varepsilon})
\]
is homeomorphic to
\[
(U,U\cap P_{1,\varepsilon}).
\]

Thus every boundary point of \(P_{1,\varepsilon}\) is locally topologically stable with respect to the parameter \(a\). Since \(\partial P_{1,\varepsilon}\) is compact, finitely many such neighborhoods \(U_1,\dots,U_N\) cover it. On the compact complement of \(\bigcup_{j=1}^N U_j\), the continuous function
\[
z\longmapsto s_{\min}(z\mathbf I_n-\mathbf A_n(1))-\varepsilon
\]
has no zeros and is therefore bounded away from \(0\). By continuity in \(a\), the corresponding sign pattern is unchanged for \(a\) sufficiently close to \(1\). Combining this with the local homeomorphisms on the \(U_j\) shows that \(a=1\) is not an exceptional parameter for \(\pi\).

Therefore there exists
\[
\widehat a_{n,\varepsilon}\in(0,1)
\]
such that the fibers
\[
P_{a,\varepsilon},
\qquad
\widehat a_{n,\varepsilon}<a\le 1,
\]
are all semialgebraically homeomorphic to \(P_{1,\varepsilon}\).

Finally, every connected component of \(P_{1,\varepsilon}\) is a connected union of equal-radius discs with centers on the real line, hence is the \(\varepsilon\)-neighborhood of a compact interval and therefore convex. In particular, every connected component of \(P_{1,\varepsilon}\) is simply connected. The same then holds for \(P_{a,\varepsilon}\) when \(a\) is sufficiently close to \(1\).
\end{proof}

\section{An algebraic boundary polynomial}

For \(z=x+iy\), define the Hermitian matrix
\[
\mathbf H_{n,\varepsilon}(x,y;a)
:=
(z\mathbf I_n-\mathbf A_n(a))^*(z\mathbf I_n-\mathbf A_n(a))
-\varepsilon^2\mathbf I_n,
\]
and its determinant
\[
F_{n,\varepsilon}(x,y;a)
:=
\det \mathbf H_{n,\varepsilon}(x,y;a).
\]

\begin{proposition}[Boundary polynomial]\label{prop:boundary-polynomial}
For every \(n\ge 2\), \(0<a<1\), and \(\varepsilon>0\), the polynomial \(F_{n,\varepsilon}(x,y;a)\) belongs to \(\R[a,x^2,y^2]\) and has total degree \(2n\) in \((x,y)\). Moreover,
\[
\partial \sigma_\varepsilon(\mathbf A_n(a))
=
\left\{
x+iy\in\C:\;
F_{n,\varepsilon}(x,y;a)=0,\;
\mathbf H_{n,\varepsilon}(x,y;a)\succeq 0
\right\}.
\]
Hence
\[
\partial \sigma_\varepsilon(\mathbf A_n(a))
\subseteq
\left\{
x+iy\in\C:\;
F_{n,\varepsilon}(x,y;a)=0
\right\}.
\]
\end{proposition}

\begin{proof}
Since the entries of \(\mathbf H_{n,\varepsilon}(x,y;a)\) are polynomial in \((a,x,y)\), the determinant \(F_{n,\varepsilon}\) is polynomial in these variables. Because \(\mathbf A_n(a)\) is real,
\[
\mathbf H_{n,\varepsilon}(x,-y;a)
=
\overline{\mathbf H_{n,\varepsilon}(x,y;a)},
\]
and therefore \(F_{n,\varepsilon}(x,-y;a)=F_{n,\varepsilon}(x,y;a)\).

Next, let
\[
S_n:=\diag(1,-1,1,-1,\dots,(-1)^{n-1}).
\]
Then \(S_n\mathbf A_n(a)S_n^{-1}=-\mathbf A_n(a)\), so
\[
S_n\mathbf H_{n,\varepsilon}(x,y;a)S_n^{-1}
=
\mathbf H_{n,\varepsilon}(-x,-y;a).
\]
Taking determinants yields
\[
F_{n,\varepsilon}(-x,-y;a)=F_{n,\varepsilon}(x,y;a),
\]
and together with the evenness in \(y\) this implies evenness in \(x\). Hence
\[
F_{n,\varepsilon}(x,y;a)\in \R[a,x^2,y^2].
\]

The degree statement is proved exactly as in the standard determinant expansion: the identity permutation contributes \((x^2+y^2)^n\), while every nonidentity permutation uses at least two off-diagonal entries and therefore lowers the total degree by at least \(2\). Thus \(\deg_{x,y}F_{n,\varepsilon}=2n\).

Now set
\[
f(z):=s_{\min}(z\mathbf I_n-\mathbf A_n(a)).
\]
Since
\[
\sigma_\varepsilon(\mathbf A_n(a))
=
\{z\in\C:\;f(z)<\varepsilon\},
\]
we always have
\[
\partial \sigma_\varepsilon(\mathbf A_n(a))
\subseteq
\{z\in\C:\;f(z)=\varepsilon\}.
\]
Conversely, if \(f(z_0)=\varepsilon\) and \(z_0\notin \partial \sigma_\varepsilon(\mathbf A_n(a))\), then \(z_0\) lies in the interior of the complement, and the resolvent norm
\[
\nu(z):=\|(z\mathbf I_n-\mathbf A_n(a))^{-1}\|
\]
attains a local maximum \(\varepsilon^{-1}\) at \(z_0\). Since \(\nu\) is subharmonic on the resolvent set, it must be locally constant; hence \(f(z)\equiv\varepsilon\) on an open set. Then \(F_{n,\varepsilon}\) vanishes on an open set, contradicting that it is a nonzero polynomial. Therefore
\[
\partial \sigma_\varepsilon(\mathbf A_n(a))
=
\{z\in\C:\;f(z)=\varepsilon\}.
\]

Finally, the eigenvalues of \(\mathbf H_{n,\varepsilon}(x,y;a)\) are \(s_j(z\mathbf I_n-\mathbf A_n(a))^2-\varepsilon^2\). Hence
\[
f(z)=\varepsilon
\iff
\lambda_{\min}\bigl(\mathbf H_{n,\varepsilon}(x,y;a)\bigr)=0
\iff
F_{n,\varepsilon}(x,y;a)=0
\ \text{and}\ 
\mathbf H_{n,\varepsilon}(x,y;a)\succeq 0.
\]
\end{proof}

\begin{remark}[Discarded branches]\label{rem:discarded-branches}
The full real algebraic curve \(F_{n,\varepsilon}(x,y;a)=0\) is generally larger than the true boundary. Indeed, the equation
\[
F_{n,\varepsilon}(x,y;a)=0
\]
merely says that \emph{some} singular value of \(z\mathbf I_n-\mathbf A_n(a)\) equals \(\varepsilon\). The true boundary is the branch where the \emph{smallest} singular value equals \(\varepsilon\).
\end{remark}

\begin{proposition}[Discarded algebraic branches lie inside]\label{prop:discarded-inside}
If
\[
F_{n,\varepsilon}(x,y;a)=0
\qquad\text{and}\qquad
\mathbf H_{n,\varepsilon}(x,y;a)\not\succeq 0,
\]
then \(x+iy\in \sigma_\varepsilon(\mathbf A_n(a))\).
\end{proposition}

\begin{proof}
The matrix \(\mathbf H_{n,\varepsilon}(x,y;a)\) is Hermitian. If it is singular and not positive semidefinite, then its smallest eigenvalue is negative. Hence
\[
s_{\min}(z\mathbf I_n-\mathbf A_n(a))^2-\varepsilon^2<0,
\]
so \(s_{\min}(z\mathbf I_n-\mathbf A_n(a))<\varepsilon\), i.e. \(z\in \sigma_\varepsilon(\mathbf A_n(a))\).
\end{proof}

\begin{proposition}[Explicit form of the Hermitian pencil]\label{prop:explicit-H}
Set
\[
\tau:=x^2+y^2-\varepsilon^2,
\qquad
c:=az+\overline z=(1+a)x-i(1-a)y,
\qquad
d:=\tau+a^2+1.
\]
Then
\[
\mathbf H_{n,\varepsilon}(x,y;a)
=
(|z|^2-\varepsilon^2)\mathbf I_n-\overline z\,\mathbf A_n(a)-z\,\mathbf A_n(a)^{\mathsf T}
+\mathbf A_n(a)^{\mathsf T}\mathbf A_n(a),
\]
and its entries are
\[
\bigl(\mathbf H_{n,\varepsilon}(x,y;a)\bigr)_{jk}
=
\begin{cases}
\tau+a^2, & j=k=1,\\[1mm]
\tau+1, & j=k=n,\\[1mm]
d, & j=k\in\{2,\dots,n-1\},\\[1mm]
-c, & k=j+1,\\[1mm]
-\overline c, & j=k+1,\\[1mm]
a, & |j-k|=2,\\[1mm]
0, & |j-k|>2.
\end{cases}
\]
\end{proposition}

\begin{proof}
The first identity is just the product expansion
\[
(\overline z\mathbf I_n-\mathbf A_n(a)^{\mathsf T})(z\mathbf I_n-\mathbf A_n(a))-\varepsilon^2\mathbf I_n.
\]
The description of the entries follows from the facts that:
\begin{itemize}
\item \(-\overline z\,\mathbf A_n(a)-z\,\mathbf A_n(a)^{\mathsf T}\) contributes the first off-diagonals \(-c\) and \(-\overline c\);
\item \(\mathbf A_n(a)^{\mathsf T}\mathbf A_n(a)\) has diagonal entries \(a^2,a^2+1,\dots,a^2+1,1\);
\item \(\mathbf A_n(a)^{\mathsf T}\mathbf A_n(a)\) has second off-diagonals equal to \(a\).
\end{itemize}
\end{proof}

\begin{corollary}[Toeplitz reduction]\label{cor:Toeplitz-reduction}
Let \(\Delta_0:=1\), and for \(m\ge 1\) define
\[
\Delta_m:=\det \Theta_m(d,c,a),
\]
where
\[
\Theta_1(d,c,a):=[d],
\qquad
\Theta_2(d,c,a):=
\begin{bmatrix}
d & -c\\
-\overline c & d
\end{bmatrix},
\]
and for \(m\ge 3\),
\[
\Theta_m(d,c,a)=
\begin{bmatrix}
d & -c & a \\
-\overline c & d & -c & \ddots \\
a & -\overline c & d & \ddots & \ddots \\
& \ddots & \ddots & \ddots & \ddots & \ddots \\
&& \ddots & \ddots & d & -c & a \\
&&& \ddots & -\overline c & d & -c \\
&&&& a & -\overline c & d
\end{bmatrix}.
\]
Then
\[
F_{n,\varepsilon}(x,y;a)
=
\Delta_n-(a^2+1)\Delta_{n-1}+a^2\Delta_{n-2}.
\]
\end{corollary}

\begin{proof}
Let \(\mathbf E_{11}\) and \(\mathbf E_{nn}\) denote the standard diagonal matrix units. By Proposition~\ref{prop:explicit-H},
\[
\mathbf H_{n,\varepsilon}(x,y;a)
=
\Theta_n(d,c,a)-\mathbf E_{11}-a^2\mathbf E_{nn}.
\]
Since the determinant is affine in each diagonal entry separately,
\[
\det\bigl(\Theta_n-\mathbf E_{11}-a^2\mathbf E_{nn}\bigr)
=
\Delta_n-\Delta_{n-1}-a^2\Delta_{n-1}+a^2\Delta_{n-2},
\]
which is exactly the stated formula.
\end{proof}

\begin{remark}[Closed form]\label{rem:quartic}
Because \(\Theta_m(d,c,a)\) is banded Toeplitz, the determinants \(\Delta_m\) satisfy a fourth-order linear recurrence in \(m\). Equivalently, \(\Delta_m\) admits a closed form built from the roots of the reciprocal quartic
\[
a\rho^4-c\rho^3+d\rho^2-\overline c\,\rho+a=0.
\]
When the quartic has repeated roots, the usual polynomial prefactors in \(m\) occur. Thus the boundary polynomial admits a closed form in \(n\), even though that form is not especially illuminating for the topology.
\end{remark}

\section{Vertical slices and the slice index}

We first record a general criterion that reduces simple connectedness to a one-dimensional statement along vertical lines.

\begin{theorem}[Vertical-slice criterion]\label{thm:vertical-slice}
Fix \(n\ge 2\), \(0<a<1\), and \(\varepsilon>0\). Assume that for every connected component \(\Omega\) of \(\sigma_\varepsilon(\mathbf A_n(a))\) and every \(x\in\R\), the set
\[
\{y\ge 0:\ x+iy\in \partial\Omega\}
\]
has at most one element. Then every connected component of \(\sigma_\varepsilon(\mathbf A_n(a))\) is simply connected.
\end{theorem}

\begin{proof}
Let \(\Omega\) be a connected component of \(\sigma_\varepsilon(\mathbf A_n(a))\). Since \(\|(z\mathbf I_n-\mathbf A_n(a))^{-1}\|\to 0\) as \(|z|\to\infty\), the set \(\Omega\) is bounded.

We claim that \(\Omega\) contains an eigenvalue of \(\mathbf A_n(a)\). If not, then \(\overline\Omega\cap \sigma(\mathbf A_n(a))=\varnothing\), so the resolvent norm
\[
\nu(z):=\|(z\mathbf I_n-\mathbf A_n(a))^{-1}\|
\]
is continuous on \(\overline\Omega\) and subharmonic in a neighborhood of \(\overline\Omega\). By Proposition~\ref{prop:boundary-polynomial},
\[
\nu(z)=\varepsilon^{-1}\quad (z\in\partial\Omega),
\]
whereas \(\nu(z)>\varepsilon^{-1}\) on \(\Omega\). This contradicts the maximum principle for subharmonic functions. Hence \(\Omega\cap \sigma(\mathbf A_n(a))\neq\varnothing\). By Lemma~\ref{lem:real-spectrum}, every eigenvalue is real, so \(\Omega\cap\R\neq\varnothing\).

Because \(\mathbf A_n(a)\) is real,
\[
\overline{\sigma_\varepsilon(\mathbf A_n(a))}=\sigma_\varepsilon(\mathbf A_n(a)),
\]
and therefore the complex conjugate
\[
\Omega^\sharp:=\{\overline z:\ z\in\Omega\}
\]
is also a connected component. Since \(\Omega\cap\R\neq\varnothing\), the components \(\Omega\) and \(\Omega^\sharp\) intersect, hence \(\Omega=\Omega^\sharp\). Thus every vertical fiber
\[
\Omega_x:=\{y\in\R:\ x+iy\in\Omega\}
\]
is symmetric with respect to \(y\mapsto -y\).

Fix \(x\in\R\). The set \(\Omega_x\) is an open bounded subset of \(\R\), symmetric under \(y\mapsto -y\). By hypothesis, its nonnegative boundary has at most one point. Therefore either \(\Omega_x=\varnothing\) or
\[
\Omega_x=(-h(x),h(x))
\]
for some \(h(x)>0\). Define
\[
I:=\{x\in\R:\ \Omega_x\neq\varnothing\}.
\]
Then \(I\) is open, and
\[
\Omega=\{x+iy:\ x\in I,\ |y|<h(x)\}.
\]
If \(I\) were disconnected, then \(\Omega\) would split accordingly, contradicting connectedness. Hence \(I\) is an interval. The deformation
\[
\Phi_t(x+iy):=x+i(1-t)y,
\qquad 0\le t\le 1,
\]
retracts \(\Omega\) onto \(I\). Thus \(\Omega\) is contractible, hence simply connected.
\end{proof}

A convenient scalar indicator for vertical slices is the negative index of the Hermitian boundary pencil.

\begin{proposition}[Slice index criterion]\label{prop:slice-index}
For fixed \(x\in\R\), define
\[
N_x(y):=\ind \mathbf H_{n,\varepsilon}(x,y;a),
\qquad y\ge 0.
\]
Then
\[
x+iy\in \sigma_\varepsilon(\mathbf A_n(a))
\iff
N_x(y)\ge 1.
\]
Equivalently,
\[
\{y\ge 0:\ x+iy\in \sigma_\varepsilon(\mathbf A_n(a))\}
=
\{y\ge 0:\ N_x(y)\ge 1\}.
\]
\end{proposition}

\begin{proof}
The eigenvalues of \(\mathbf H_{n,\varepsilon}(x,y;a)\) are
\[
s_j((x+iy)\mathbf I_n-\mathbf A_n(a))^2-\varepsilon^2,
\qquad j=1,\dots,n.
\]
Hence \(\mathbf H_{n,\varepsilon}(x,y;a)\) has a negative eigenvalue if and only if some singular value of \(((x+iy)\mathbf I_n-\mathbf A_n(a))\) is strictly smaller than \(\varepsilon\), which is equivalent to \(x+iy\in \sigma_\varepsilon(\mathbf A_n(a))\).
\end{proof}

\begin{corollary}[Index monotonicity implies simple connectedness]\label{cor:index-monotone}
Assume that for every \(x\in\R\), the function \(y\mapsto N_x(y)\) is nonincreasing on \([0,\infty)\). Then every connected component of \(\sigma_\varepsilon(\mathbf A_n(a))\) is simply connected.
\end{corollary}

\begin{proof}
By Proposition~\ref{prop:slice-index},
\[
I_x^+:=\{y\ge 0:\ x+iy\in \sigma_\varepsilon(\mathbf A_n(a))\}
=
\{y\ge 0:\ N_x(y)\ge 1\}.
\]
Since \(N_x\) is integer-valued and nonincreasing, \(I_x^+\) is an initial segment of \([0,\infty)\). Because the pseudospectrum is open, either \(I_x^+=\varnothing\) or \(I_x^+=[0,h_x)\) for some \(h_x>0\). By conjugation symmetry, the full slice is either empty or
\[
(-h_x,h_x).
\]
Thus the hypothesis of Theorem~\ref{thm:vertical-slice} is satisfied.
\end{proof}

\begin{remark}
Corollary~\ref{cor:index-monotone} is the natural point at which discrete oscillation theory and matrix Pr\"ufer methods enter; compare \cite{Hilscher2012,KratzHilscher2013,SchulzBaldes2007}. The final section below records a different, explicitly conditional route: rather than proving global slice monotonicity directly, it isolates a local topological mechanism that would rule out hole formation on the remaining compact parameter strip.
\end{remark}

\section{A conditional componentwise topological argument on the remaining strip}

The purpose of this section is to isolate the local topological step that would exclude hole formation on the remaining compact parameter strip once the requisite analytic input has been established elsewhere. Accordingly, the statements explicitly labelled as hypotheses below are external inputs rather than results proved in this note.

For fixed \(n\ge 2\) and \(\varepsilon>0\), set
\[
P_{a,\varepsilon}:=\sigma_\varepsilon(\mathbf A_n(a)),
\qquad 0<a\le 1.
\]

\begin{definition}[Topology-changing boundary point]\label{def:topology-changing}
Fix \(a_*\in(0,1]\). A point \(p\in \partial P_{a_*,\varepsilon}\subset\R^2\) is called \emph{topology-changing} if for every neighborhood \(U\subset\R^2\) of \(p\) and every \(\delta>0\) there exist \(a_-,a_+\in (a_* -\delta,a_*+\delta)\cap(0,1]\) such that the pairs
\[
\bigl(U,U\cap P_{a_-,\varepsilon}\bigr)
\qquad\text{and}\qquad
\bigl(U,U\cap P_{a_+,\varepsilon}\bigr)
\]
are not homeomorphic.
\end{definition}

\begin{hypothesis}[Chart coverage on the remaining strip]\label{hyp:standing-input}
Fix \(n\ge 2\) and \(\varepsilon>0\). Let \(a_-\) be as in Proposition~\ref{prop:low-a}, assume that \(\varepsilon\) is noncritical at the normal endpoint, and let
\[
\widehat a_{n,\varepsilon}
\]
be as in Proposition~\ref{prop:near-normal}. Assume that on the compact strip
\[
I:=[a_-,\widehat a_{n,\varepsilon}],
\]
every boundary point of \(P_{a,\varepsilon}\) is represented in the chart introduced below, namely by an equation \(\Phi_a(u,v)=0\) with
\[
\mathbf N_a(t(u,v),\beta(u,v))\succ 0.
\]
\end{hypothesis}

\begin{remark}
Proposition~\ref{prop:low-a} and Proposition~\ref{prop:near-normal} treat the two endpoint regimes. Hypothesis~\ref{hyp:standing-input} is only the remaining chart-coverage input on the compact strip \(I\).
\end{remark}

\subsection*{A boundary chart on the remaining strip}

Fix \(a\in I\), and put
\[
\tau:=1-a>0.
\]
Introduce the upper-half-plane variables
\[
u:=x^2,\qquad v:=y^2,
\]
and define
\[
\beta(u,v):=\frac{u+v-\varepsilon^2+\tau^2}{a},
\qquad
t(u,v):=\frac{(1+a)^2u+\tau^2v}{4a^2}.
\]
Let
\[
\mathbf T_n:=\tridiag(1,0,1),
\qquad
\mathbf E_{nn}:=(\delta_{jn}\delta_{kn})_{j,k=1}^n,
\qquad
\mathbf e_1:=(1,0,\dots,0)^{\mathsf T}.
\]
On the chart under consideration one arrives at the \(n\times n\) Hermitian matrix
\[
\mathbf N_a(t,\beta)
:=
\mathbf T_n^2-2\sqrt t\,\mathbf T_n+\beta\mathbf I_n+\tau\mathbf E_{nn}
=
(\mathbf T_n-\sqrt t\,\mathbf I_n)^2+(\beta-t)\mathbf I_n+\tau\mathbf E_{nn}.
\]
Along this chart one has \(\mathbf N_a(t,\beta)\succ 0\), so the scalar function
\[
m_n(t,\beta;a)
:=
\mathbf e_1^{\mathsf T}\mathbf N_a(t,\beta)^{-1}\mathbf e_1
\]
is well defined. The boundary equation in upper-half-plane variables is then
\[
\Phi_a(u,v)
:=
m_n(t(u,v),\beta(u,v);a)-\frac{a}{\tau}
=
0.
\]
We also set
\[
\widetilde\Phi_a(x,v):=\Phi_a(x^2,v),
\qquad v\ge 0.
\]

\begin{remark}\label{rem:rational}
By the adjugate formula, the entries of \(\mathbf N_a(t,\beta)^{-1}\) are rational functions of \(\sqrt t\), \(\beta\), and \(a\). Hence \(m_n\), \(\Phi_a\), and \(\widetilde\Phi_a\) are real-analytic and semialgebraic on the chart.
\end{remark}

\begin{hypothesis}[Monotonicity input on the chart]\label{hyp:monotone-m}
For every \(a\in I\) and every point of the chart,
\[
\partial_t m_n(t,\beta;a)>0,
\qquad
\partial_\beta m_n(t,\beta;a)<0.
\]
Moreover,
\[
\partial_\beta m_n(t,\beta;a)
=
-\mathbf e_1^{\mathsf T}\mathbf N_a(t,\beta)^{-2}\mathbf e_1<0.
\]
\end{hypothesis}

\begin{lemma}[Derivative identity]\label{lem:derivative-identity}
For every \(a\in I\),
\[
\partial_u\Phi_a
=
\frac1a\,\partial_\beta m_n
+
\frac{(1+a)^2}{4a^2}\,\partial_t m_n,
\]
\[
\partial_v\Phi_a
=
\frac1a\,\partial_\beta m_n
+
\frac{\tau^2}{4a^2}\,\partial_t m_n.
\]
Hence
\[
\partial_u\Phi_a-\partial_v\Phi_a
=
\frac1a\,\partial_t m_n
>0.
\]
\end{lemma}

\begin{proof}
This is a direct chain-rule computation using
\[
\partial_u\beta=\partial_v\beta=\frac1a,
\qquad
\partial_u t=\frac{(1+a)^2}{4a^2},
\qquad
\partial_v t=\frac{\tau^2}{4a^2}.
\]
Subtracting the two identities gives
\[
\partial_u\Phi_a-\partial_v\Phi_a
=
\frac{(1+a)^2-\tau^2}{4a^2}\,\partial_t m_n
=
\frac{4a}{4a^2}\,\partial_t m_n
=
\frac1a\,\partial_t m_n>0.
\]
\end{proof}

\begin{lemma}[No off-axis critical points]\label{lem:no-off-axis-critical}
Fix \(a\in I\). Let \(x\neq 0\) and \(v>0\). If
\[
\widetilde\Phi_a(x,v)=0,
\]
then
\[
\nabla \widetilde\Phi_a(x,v)\neq 0.
\]
\end{lemma}

\begin{proof}
We have
\[
\partial_x\widetilde\Phi_a(x,v)=2x\,\partial_u\Phi_a(x^2,v),
\qquad
\partial_v\widetilde\Phi_a(x,v)=\partial_v\Phi_a(x^2,v).
\]
If both vanished, then \(x\neq 0\) would force
\[
\partial_u\Phi_a(x^2,v)=0,
\qquad
\partial_v\Phi_a(x^2,v)=0,
\]
contradicting Lemma~\ref{lem:derivative-identity}.
\end{proof}

\begin{lemma}[Real-axis tangencies]\label{lem:real-axis-tangency}
Fix \(a\in I\), and let \(x_0\neq 0\). Suppose
\[
\widetilde\Phi_a(x_0,0)=0,
\qquad
\partial_x\widetilde\Phi_a(x_0,0)=0.
\]
Then
\[
\partial_v\widetilde\Phi_a(x_0,0)<0.
\]
Consequently, near \((x_0,0)\), the zero set of \(\widetilde\Phi_a\) is the graph
\[
v=\psi(x),
\qquad
\psi(x_0)=0,
\qquad
\psi'(x_0)=0.
\]
Equivalently, in \((x,y)\)-coordinates the boundary is locally of the form
\[
y^2=\psi(x).
\]
\end{lemma}

\begin{proof}
Since \(x_0\neq 0\),
\[
0=\partial_x\widetilde\Phi_a(x_0,0)=2x_0\,\partial_u\Phi_a(x_0^2,0),
\]
hence \(\partial_u\Phi_a(x_0^2,0)=0\). By Lemma~\ref{lem:derivative-identity},
\[
0<\partial_u\Phi_a(x_0^2,0)-\partial_v\Phi_a(x_0^2,0),
\]
so \(\partial_v\Phi_a(x_0^2,0)<0\), i.e.
\[
\partial_v\widetilde\Phi_a(x_0,0)<0.
\]
The implicit-function theorem now gives the stated graph \(v=\psi(x)\), and differentiating the identity \(\widetilde\Phi_a(x,\psi(x))=0\) at \(x=x_0\) yields \(\psi'(x_0)=0\).
\end{proof}

\begin{lemma}[Imaginary-axis tangencies]\label{lem:imag-axis-tangency}
Fix \(a\in I\), and let \(v_0>0\). Suppose
\[
\widetilde\Phi_a(0,v_0)=0,
\qquad
\partial_v\widetilde\Phi_a(0,v_0)=0.
\]
Then
\[
\partial_u\Phi_a(0,v_0)>0.
\]
Consequently, near \((0,v_0)\) in \((u,v)\)-coordinates the boundary is a graph
\[
u=\phi(v),
\qquad
\phi(v_0)=0,
\qquad
\phi'(v_0)=0.
\]
If \(y_0:=\sqrt{v_0}\), then near \((0,\pm y_0)\) in \((x,y)\)-coordinates the boundary is locally of the form
\[
x^2=\chi_\pm(y),
\qquad
\chi_\pm(\pm y_0)=0,
\qquad
\chi_\pm'(\pm y_0)=0.
\]
\end{lemma}

\begin{proof}
At \(x=0\), one automatically has \(\partial_x\widetilde\Phi_a(0,v)=0\). The hypothesis gives
\[
0=\partial_v\widetilde\Phi_a(0,v_0)=\partial_v\Phi_a(0,v_0).
\]
By Lemma~\ref{lem:derivative-identity},
\[
\partial_u\Phi_a(0,v_0)>\partial_v\Phi_a(0,v_0)=0.
\]
Hence \(\partial_u\Phi_a(0,v_0)>0\), and the implicit-function theorem yields the graph \(u=\phi(v)\) with \(\phi(v_0)=0\). Differentiating \(\Phi_a(\phi(v),v)=0\) at \(v=v_0\) gives \(\phi'(v_0)=0\). Passing from \(v=y^2\) back to \(y\) yields the two local models \(x^2=\chi_\pm(y)\) near \((0,\pm y_0)\).
\end{proof}

\begin{remark}[Geometric content of the axis-tangency lemmas]\label{rem:fold-geometry}
Lemmas~\ref{lem:real-axis-tangency} and \ref{lem:imag-axis-tangency} show that, away from the origin, axis singularities detected by the chart are fold-type in the sense that the zero set is locally given by \(y^2=\psi(x)\) or \(x^2=\chi(y)\). These lemmas describe only the local geometry of the zero set. They do not by themselves determine which side of the local graph belongs to \(P_{a,\varepsilon}\).
\end{remark}

\begin{proposition}[The origin is not topology-changing]\label{prop:origin-not-topology-changing}
Assume Hypothesis~\ref{hyp:standing-input} and Hypothesis~\ref{hyp:monotone-m}. Then for every
\[
a_*\in I,
\]
the origin is not a topology-changing boundary point of the family
\[
\{P_{a,\varepsilon}\}_{a\in I}.
\]
\end{proposition}

\begin{proof}
We split into the parity of \(n\).

If \(n\) is odd, then by Lemma~\ref{lem:real-spectrum},
\[
0=2\sqrt a\cos\frac{(n+1)\pi}{2(n+1)}\in \sigma(\mathbf A_n(a))
\]
for every \(a\in I\). Hence
\[
0\in \sigma(\mathbf A_n(a))\subset P_{a,\varepsilon},
\]
and since \(P_{a,\varepsilon}\) is open, the origin is never a boundary point at all. In particular, it cannot be topology-changing.

Now assume that \(n\) is even, and let \(a_*\in I\) be such that
\[
0\in \partial P_{a_*,\varepsilon}.
\]
By Hypothesis~\ref{hyp:standing-input}, the origin is represented in the chart, so
\[
\Phi_{a_*}(0,0)=0.
\]
By Lemma~\ref{lem:derivative-identity},
\[
\partial_u\Phi_{a_*}(0,0)-\partial_v\Phi_{a_*}(0,0)>0,
\]
hence
\[
\nabla \Phi_{a_*}(0,0)\neq 0.
\]

If \(\partial_v\Phi_{a_*}(0,0)=0\), then \(\partial_u\Phi_{a_*}(0,0)\neq 0\), and the same argument below applies with the roles of \(u\) and \(v\) interchanged. Thus it suffices to treat the case
\[
\partial_v\Phi_{a_*}(0,0)\neq 0.
\]
By the implicit-function theorem with parameter, there exist numbers \(\rho>0\), \(\delta>0\), an interval
\[
J:=(a_*-\delta,a_*+\delta)\cap I,
\]
and a real-analytic function
\[
\psi:[0,\rho^2]\times J\to[0,\rho^2]
\]
such that
\[
\psi(0,a_*)=0
\]
and, for every \(a\in J\),
\[
[0,\rho^2]^2\cap \{\Phi_a(u,v)=0\}
=
\{(u,v)\in [0,\rho^2]^2:\ v=\psi(u,a)\}.
\]
After shrinking \(\rho\) and \(\delta\) if necessary, we may also assume that
\[
\partial_v\Phi_a(u,\psi(u,a))\neq 0
\]
for all \(u\in[0,\rho^2]\) and \(a\in J\). Thus in the first quadrant the zero set is a single graph depending continuously on \(a\), and the local side decomposition is constant in \(a\).

Let
\[
Q_+:=[0,\rho^2]^2,
\qquad
U:=[-\rho,\rho]^2.
\]
Define
\[
\widehat P_a
:=
\{(u,v)\in Q_+:\ \sqrt u+i\sqrt v\in P_{a,\varepsilon}\}.
\]
Because \(P_{a,\varepsilon}\) is invariant under \(x\mapsto -x\) and \(y\mapsto -y\), this first-quadrant set determines the full local pair near the origin. Since the boundary in \(Q_+\) is the graph \(v=\psi(u,a)\), the set \(\widehat P_a\) is, for every \(a\in J\), exactly one of the two graph sides
\[
\{(u,v)\in Q_+:\ v>\psi(u,a)\}
\quad\text{or}\quad
\{(u,v)\in Q_+:\ v<\psi(u,a)\}.
\]
The choice is independent of \(a\in J\), because the sign of \(\Phi_a\) at any fixed point not on the graph varies continuously with \(a\) and cannot change without crossing the graph.

For each \(u\in[0,\rho^2]\) and \(a\in J\), let
\[
\eta_{u,a}:[0,\rho^2]\to[0,\rho^2]
\]
be the piecewise linear homeomorphism fixing \(0\) and \(\rho^2\) and sending
\[
\psi(u,a)\mapsto \psi(u,a_*).
\]
Then
\[
h_a(u,v):=(u,\eta_{u,a}(v))
\]
is a homeomorphism of \(Q_+\) onto itself, and by construction
\[
h_a(\widehat P_a)=\widehat P_{a_*}.
\]

Now define \(H_a:U\to U\) by
\[
H_a(x,y)
:=
\left(
x,\;
\operatorname{sgn}(y)\sqrt{\eta_{x^2,a}(y^2)}
\right),
\qquad \operatorname{sgn}(0):=0.
\]
Since each \(\eta_{x^2,a}\) is a homeomorphism of \([0,\rho^2]\), the map \(H_a\) is a homeomorphism of \(U\), with inverse
\[
H_a^{-1}(x,y)
=
\left(
x,\;
\operatorname{sgn}(y)\sqrt{\eta_{x^2,a}^{-1}(y^2)}
\right).
\]
By reflection symmetry of \(P_{a,\varepsilon}\) in both coordinate axes and by the definition of \(\widehat P_a\), we obtain
\[
H_a\bigl(U\cap P_{a,\varepsilon}\bigr)=U\cap P_{a_*,\varepsilon}.
\]
Thus the local pairs
\[
\bigl(U,U\cap P_{a,\varepsilon}\bigr),
\qquad a\in J,
\]
are mutually homeomorphic. Therefore the origin is not topology-changing.
\end{proof}

\begin{definition}[Minimum-type fold models]\label{def:min-fold}
Let
\[
Q:=[-1,1]^2,
\qquad
\Omega_{\mathrm{re}}:=\{(\xi,\eta)\in Q:\ \eta^2<\xi\},
\qquad
\Omega_{\mathrm{im}}:=\{(\xi,\eta)\in Q:\ \xi^2<\eta\}.
\]
A local pair \((U,U\cap\Omega)\) is called a \emph{real-axis minimum-type fold} if it is semialgebraically homeomorphic to \((Q,\Omega_{\mathrm{re}})\). It is called an \emph{imaginary-axis minimum-type fold} if it is semialgebraically homeomorphic to \((Q,\Omega_{\mathrm{im}})\).
\end{definition}

\begin{proposition}[Minimum-type fold models are hole-free]\label{prop:min-fold-hole-free}
The sets \(\Omega_{\mathrm{re}}\) and \(\Omega_{\mathrm{im}}\) are contractible. Their complements in \(Q\) are contractible as well. In particular, neither model traps a bounded complementary component away from \(\partial Q\).
\end{proposition}

\begin{proof}
For \(\Omega_{\mathrm{re}}\), the homotopy
\[
H_t(\xi,\eta):=\bigl((1-t)\xi+t/2,(1-t)\eta\bigr),
\qquad 0\le t\le 1,
\]
remains in \(\Omega_{\mathrm{re}}\) because
\[
((1-t)\eta)^2<(1-t)^2\xi<(1-t)\xi+t/2.
\]
Thus \(\Omega_{\mathrm{re}}\) contracts to \((1/2,0)\). Its complement in \(Q\) is \(\{(\xi,\eta)\in Q:\ \xi\le \eta^2\}\), and the homotopy
\[
K_t(\xi,\eta):=\bigl((1-t)\xi-t,\eta\bigr),
\qquad 0\le t\le 1,
\]
stays in the complement because \((1-t)\xi-t\le \xi\le \eta^2\). Hence \(Q\setminus\Omega_{\mathrm{re}}\) contracts to the left edge \(\{-1\}\times[-1,1]\).

For \(\Omega_{\mathrm{im}}\), the homotopy
\[
L_t(\xi,\eta):=\bigl((1-t)\xi,(1-t)\eta+t/2\bigr),
\qquad 0\le t\le 1,
\]
remains in \(\Omega_{\mathrm{im}}\) because
\[
((1-t)\xi)^2<(1-t)^2\eta<(1-t)\eta+t/2.
\]
Thus \(\Omega_{\mathrm{im}}\) contracts to \((0,1/2)\). Its complement is \(\{(\xi,\eta)\in Q:\ \eta\le \xi^2\}\), and the homotopy
\[
M_t(\xi,\eta):=\bigl(\xi,(1-t)\eta-t\bigr),
\qquad 0\le t\le 1,
\]
stays in the complement because \((1-t)\eta-t\le \eta\le \xi^2\). Hence \(Q\setminus\Omega_{\mathrm{im}}\) contracts to the bottom edge \([-1,1]\times\{-1\}\).
\end{proof}

\begin{hypothesis}[Minimum-type character of topology-changing axis tangencies]\label{hyp:min-fold}
Whenever \(p\in\partial P_{a_*,\varepsilon}\) with \(a_*\in I\) is a topology-changing point lying on the real or imaginary axis away from the origin, and \(p\) is detected by Lemma~\ref{lem:real-axis-tangency} or Lemma~\ref{lem:imag-axis-tangency}, there exists a neighborhood \(U\) of \(p\) such that the pair \((U,U\cap P_{a_*,\varepsilon})\) is a minimum-type fold model in the sense of Definition~\ref{def:min-fold}.
\end{hypothesis}

\begin{remark}
Hypothesis~\ref{hyp:min-fold} is precisely the missing side-information from Remark~\ref{rem:fold-geometry}. It rules out saddle-type local models at axis tangencies relevant to topology change.
\end{remark}

\subsection*{Topological preliminaries}

\begin{lemma}[Planar criterion]\label{lem:planar-criterion}
A connected open set \(\Omega\subset\R^2\) is simply connected if and only if \(\R^2\setminus\Omega\) has no bounded connected component.
\end{lemma}

\begin{proof}
This is standard planar topology: after one-point compactification, \(\Omega\) is simply connected if and only if its complement in \(S^2\) is connected. Since the component containing \(\infty\) is always present, this is equivalent to saying that \(\R^2\setminus\Omega\) has no bounded connected component.
\end{proof}

\begin{lemma}[Uniform boundedness and semialgebraicity]\label{lem:semialg}
Fix \(\varepsilon>0\) and let \(J\subset (0,1]\) be compact. Then
\[
\mathcal X_J
:=
\{(x,y,a)\in \R^2\times J:\ x+iy\in P_{a,\varepsilon}\}
\]
is bounded and semialgebraic.
\end{lemma}

\begin{proof}
Write
\[
\mathbf A_n(a)=\mathbf U+a\mathbf L,
\]
where \(\mathbf U\) is the upper shift and \(\mathbf L\) is the lower shift. Since \(\|\mathbf U\|_2=\|\mathbf L\|_2=1\) and \(0<a\le 1\),
\[
\|\mathbf A_n(a)\|_2\le 2.
\]
If \(z\in P_{a,\varepsilon}\), then
\[
s_{\min}(z\mathbf I_n-\mathbf A_n(a))<\varepsilon,
\]
but also
\[
s_{\min}(z\mathbf I_n-\mathbf A_n(a))
\ge |z|-\|\mathbf A_n(a)\|_2
\ge |z|-2.
\]
Hence \(|z|<2+\varepsilon\), uniformly for \(a\in J\). This proves boundedness.

For semialgebraicity, write \(z=x+iy\) and \(\boldsymbol\xi=\mathbf u+i\mathbf v\) with \(\mathbf u,\mathbf v\in\R^n\). Then
\[
(x,y,a)\in \mathcal X_J
\]
if and only if there exist \(\mathbf u,\mathbf v\in\R^n\) such that
\[
\|\mathbf u\|^2+\|\mathbf v\|^2=1
\]
and
\[
\|((x\mathbf I_n-\mathbf A_n(a))\mathbf u-y\mathbf v)\|^2
+
\|(y\mathbf u+(x\mathbf I_n-\mathbf A_n(a))\mathbf v)\|^2
<\varepsilon^2.
\]
These are polynomial equalities and inequalities in real variables, so \(\mathcal X_J\) is semialgebraic by the Tarski--Seidenberg theorem.
\end{proof}

\begin{lemma}[Local triviality away from critical events]\label{lem:local-trivial}
Fix \(a_*\in I\), and let \(p\in \partial P_{a_*,\varepsilon}\). Assume that \(p\) is represented in the chart and satisfies one of the following:
\begin{enumerate}
\item \(p=(x_0,y_0)\) with \(x_0\neq 0\) and \(y_0\neq 0\);
\item \(p=(x_0,0)\) with \(x_0\neq 0\) and \(\partial_x\widetilde\Phi_{a_*}(x_0,0)\neq 0\);
\item \(p=(0,\pm y_0)\) with \(y_0>0\) and \(\partial_v\widetilde\Phi_{a_*}(0,y_0^2)\neq 0\).
\end{enumerate}
Then there exist a neighborhood \(U\subset\R^2\) of \(p\) and \(\delta>0\) such that the pairs
\[
\bigl(U,U\cap P_{a,\varepsilon}\bigr),
\qquad
|a-a_*|<\delta,
\]
are mutually homeomorphic.
\end{lemma}

\begin{proof}
In case \((1)\), Lemma~\ref{lem:no-off-axis-critical} and the implicit-function theorem show that the boundary is a smooth embedded arc depending continuously on \(a\). In case \((2)\), the implicit-function theorem writes the boundary locally as \(x=\chi(v,a)\). In case \((3)\), the implicit-function theorem writes the boundary locally as \(v=\nu(x,a)\), equivalently as an embedded arc in \((x,y)\)-coordinates. In each case, the local side decomposition is constant for \(a\) near \(a_*\).
\end{proof}

\begin{theorem}[Conditional componentwise simple connectedness]\label{thm:main}
Fix \(n\ge 2\) and \(\varepsilon>0\), and assume Hypotheses~\ref{hyp:standing-input}, \ref{hyp:monotone-m}, and \ref{hyp:min-fold} for these parameters. Then every connected component of
\[
P_{a,\varepsilon}
=
\sigma_\varepsilon(\mathbf A_n(a)),
\qquad 0<a<1,
\]
is simply connected.

Consequently, if
\[
0<\alpha<\beta,
\]
then
\[
\tridiag(\alpha,0,\beta)=\beta\,\mathbf A_n(\eta),
\qquad
\eta:=\alpha/\beta\in(0,1),
\]
and therefore
\[
\sigma_\varepsilon(\tridiag(\alpha,0,\beta))
=
\beta\,\sigma_{\varepsilon/\beta}(\mathbf A_n(\eta)).
\]
Hence the same conclusion holds for
\[
\sigma_\varepsilon(\tridiag(\alpha,0,\beta))
\]
after normalizing to \(\eta=\alpha/\beta\) and \(\varepsilon/\beta\) and assuming the corresponding hypotheses for the normalized family.
\end{theorem}

\begin{proof}[Proof sketch under the stated hypotheses]
By Proposition~\ref{prop:low-a}, the statement is already known in the low-\(a\) regime \(0<a<a_-\). By Proposition~\ref{prop:near-normal}, it is also known in the near-normal regime
\[
\widehat a_{n,\varepsilon}<a\le 1.
\]
Thus only the compact strip
\[
I=[a_-,\widehat a_{n,\varepsilon}]
\]
remains.

Choose \(a_0\in(0,a_-)\) and consider
\[
\mathcal X
:=
\{(x,y,a)\in \R^2\times [a_0,\widehat a_{n,\varepsilon}]:\ x+iy\in P_{a,\varepsilon}\}.
\]
By Lemma~\ref{lem:semialg}, \(\mathcal X\) is bounded and semialgebraic. Hence Hardt's semialgebraic triviality theorem (see, e.g., \cite{BCR}) implies that the projection
\[
\pi:\mathcal X\to [a_0,\widehat a_{n,\varepsilon}],\qquad (x,y,a)\mapsto a,
\]
has only finitely many exceptional parameter values. On each component of the complement of that finite set, the fibers \(P_{a,\varepsilon}\) are semialgebraically homeomorphic.

Suppose, for contradiction, that the theorem fails. Since all fibers with \(a<a_-\) have simply connected components, there exists a first exceptional value
\[
a_*\in I
\]
such that for \(a>a_*\) sufficiently close to \(a_*\), some connected component of \(P_{a,\varepsilon}\) is not simply connected, whereas for \(a<a_*\) sufficiently close to \(a_*\), every connected component is simply connected.

By Lemma~\ref{lem:planar-criterion}, this means that when \(a\) crosses \(a_*\) from left to right, some connected component acquires a new bounded complementary component, that is, a new hole.

We now classify the possible local topology-changing points on \(\partial P_{a_*,\varepsilon}\).

\begin{itemize}
\item By Proposition~\ref{prop:origin-not-topology-changing}, the origin is not a topology-changing point.

\item Let \(p\in\partial P_{a_*,\varepsilon}\) be topology-changing and different from the origin. If \(p\) is off the coordinate axes, then Lemma~\ref{lem:no-off-axis-critical} shows that it is regular; hence Lemma~\ref{lem:local-trivial} gives local topological constancy.

\item If \(p=(x_0,0)\) with \(x_0\neq 0\), then either \(\partial_x\widetilde\Phi_{a_*}(x_0,0)\neq 0\), in which case Lemma~\ref{lem:local-trivial} applies, or \(\partial_x\widetilde\Phi_{a_*}(x_0,0)=0\), in which case Lemma~\ref{lem:real-axis-tangency} identifies a fold-type axis tangency. By Hypothesis~\ref{hyp:min-fold}, the corresponding local pair is of minimum type, and Proposition~\ref{prop:min-fold-hole-free} shows that such a neighborhood cannot trap a bounded complementary component.

\item If \(p=(0,\pm y_0)\) with \(y_0>0\), then either \(\partial_v\widetilde\Phi_{a_*}(0,y_0^2)\neq 0\), in which case Lemma~\ref{lem:local-trivial} applies, or \(\partial_v\widetilde\Phi_{a_*}(0,y_0^2)=0\), in which case Lemma~\ref{lem:imag-axis-tangency} identifies a fold-type axis tangency. Again Hypothesis~\ref{hyp:min-fold} and Proposition~\ref{prop:min-fold-hole-free} exclude local hole creation.
\end{itemize}

Thus every topology-changing point at \(a=a_*\) is either locally trivial or else modeled by a minimum-type fold neighborhood, and the latter are hole-free by Proposition~\ref{prop:min-fold-hole-free}. Since the family is locally trivial away from finitely many such neighborhoods, crossing \(a=a_*\) cannot create a new bounded complementary component of any connected component. This contradicts the choice of \(a_*\).

Therefore no such bad parameter exists on \(I\). Combining this with Proposition~\ref{prop:low-a} and Proposition~\ref{prop:near-normal} proves the theorem.
\end{proof}

\begin{remark}
Theorem~\ref{thm:main} is explicitly conditional. In the present reduction, the low-\(a\) regime is handled by Proposition~\ref{prop:low-a}, and the noncritical near-normal regime by Proposition~\ref{prop:near-normal}. To obtain an unconditional proof one would still need to establish Hypotheses~\ref{hyp:standing-input}, \ref{hyp:monotone-m}, and \ref{hyp:min-fold}, and to treat the critical normal-endpoint radii excluded by Definition~\ref{def:noncritical}. The content of the theorem is componentwise rather than global: disconnected pseudospectral components may coexist, but under the stated assumptions no hole can be born inside any one of them.
\end{remark}

\begin{thebibliography}{99}

\bibitem{BCR}
J.~Bochnak, M.~Coste, and M.-F.~Roy,
\newblock \emph{Real Algebraic Geometry},
\newblock Ergebnisse der Mathematik und ihrer Grenzgebiete (3), Vol.~36,
Springer, Berlin, 1998.

\bibitem{Hilscher2012}
R.~{\v S}imon Hilscher,
\newblock Oscillation theorems for discrete symplectic systems with nonlinear dependence in spectral parameter,
\newblock \emph{Linear Algebra Appl.} \textbf{437} (2012), no.~12, 2922--2960.

\bibitem{KratzHilscher2013}
W.~Kratz and R.~{\v S}imon Hilscher,
\newblock A generalized index theorem for monotone matrix-valued functions with applications to discrete oscillation theory,
\newblock \emph{SIAM J. Matrix Anal. Appl.} \textbf{34} (2013), no.~1, 228--243.

\bibitem{SchulzBaldes2007}
H.~Schulz-Baldes,
\newblock Rotation numbers for Jacobi matrices with matrix entries,
\newblock \emph{Math. Phys. Electron. J.} \textbf{13} (2007), Paper No.~5, 40~pp.

\bibitem{TrefethenEmbree2005}
L.~N.~Trefethen and M.~Embree,
\newblock \emph{Spectra and Pseudospectra: The Behavior of Nonnormal Matrices and Operators},
\newblock Princeton University Press, Princeton, NJ, 2005.

\end{thebibliography}

\end{document}
